\begin{mistake}{2026-04-12}{11}

\begin{problemcontent}
求与
\[
\alpha_1=(1,2,3,-1)^{\mathrm T},\quad
\alpha_2=(0,1,1,2)^{\mathrm T},\quad
\alpha_3=(2,1,3,0)^{\mathrm T}
\]
都正交的单位向量。

\end{problemcontent}

\begin{solutioncontent}
设所求向量为
\[
x=(x_1,x_2,x_3,x_4)^{\mathrm T}.
\]

由“与 \(\alpha_1,\alpha_2,\alpha_3\) 都正交”得
\[
\alpha_i^{\mathrm T}x=0\quad (i=1,2,3).
\]

故满足齐次线性方程组
\[
\begin{bmatrix}
1&2&3&-1\\
0&1&1&2\\
2&1&3&0
\end{bmatrix}
x=0.
\]

对系数矩阵作行变换：
\[
\begin{aligned}
\begin{bmatrix}
1&2&3&-1\\
0&1&1&2\\
2&1&3&0
\end{bmatrix}
&\xrightarrow{R_3-2R_1}
\begin{bmatrix}
1&2&3&-1\\
0&1&1&2\\
0&-3&-3&2
\end{bmatrix}\\
&\xrightarrow{R_1-2R_2,\,R_3+3R_2}
\begin{bmatrix}
1&0&1&-5\\
0&1&1&2\\
0&0&0&8
\end{bmatrix}\\
&\xrightarrow{R_3/8,\,R_1+5R_3,\,R_2-2R_3}
\begin{bmatrix}
1&0&1&0\\
0&1&1&0\\
0&0&0&1
\end{bmatrix}.
\end{aligned}
\]

对应方程组为
\[
\begin{cases}
x_1+x_3=0,\\
x_2+x_3=0,\\
x_4=0.
\end{cases}
\]

令自由变量 \(x_3=t\)，则
\[
x_1=-t,\qquad x_2=-t,\qquad x_4=0.
\]

故通解为
\[
x=t(-1,-1,1,0)^{\mathrm T}.
\]

取基础解向量
\[
\eta=(-1,-1,1,0)^{\mathrm T}.
\]

其长度为
\[
\|\eta\|=\sqrt{(-1)^2+(-1)^2+1^2}=\sqrt{3}.
\]

所以所求单位向量为
\[
\pm\frac{1}{\sqrt{3}}(-1,-1,1,0)^{\mathrm T}
=
\pm\left(-\frac{1}{\sqrt{3}},-\frac{1}{\sqrt{3}},\frac{1}{\sqrt{3}},0\right)^{\mathrm T}.
\]

\end{solutioncontent}

\mistakeknowledge{
\begin{itemize}
  \item \textbf{核心公式/定理}：与若干向量都正交的问题，本质是求由这些向量作行组成的矩阵的零空间，再把非零解单位化。
  \item \textbf{易错点}：求出方向向量后还没结束，题目要求“单位向量”，必须再除以模长；另外答案应写成 \(\pm\) 两个方向。
  \item \textbf{关联知识}：若齐次方程组有一个自由变量，则解空间是一维；任一非零解都可作为基础解向量，再通过单位化得到标准正交方向。
\end{itemize}}

\end{mistake}
